\chapter{External Galaxies}

\section{Introduction}

The study of \textbf{external galaxies} marks a fundamental step beyond the description of the Milky Way, extending astrophysical analysis to systems distributed across cosmological distances. Historically, the understanding of galaxies evolved gradually through observational advances and conceptual breakthroughs.

At the end of the 18th century, the first systematic catalogues of nebulae were compiled. The \textbf{Messier catalogue} represented an early attempt to classify diffuse objects in the sky, followed by more extensive surveys such as the Herschel catalogues, which contained thousands of nebulae. Later developments led to the \textbf{New General Catalogue (NGC)} and the \textbf{Index Catalogue (IC)}, which remain widely used today.

A crucial turning point occurred in the early 20th century with the so-called \textbf{Great Debate}, concerning the nature of spiral nebulae. The question was whether these objects were part of the Milky Way or independent systems. This issue was resolved by \textbf{Hubble}, who measured distances to Cepheid variables in external systems, demonstrating that these nebulae are in fact galaxies far beyond our own.

\begin{observationblock}[Limits of early models]
Early models of the Milky Way, such as those proposed by Kapteyn, were affected by the neglect of interstellar extinction. This led to an underestimation of the true size of the Galaxy and an incorrect placement of the Sun near its center.
\end{observationblock}

\section{Morphology}

The \textbf{morphological classification} of galaxies provides a systematic framework to describe their structure and appearance. The most widely used scheme is the \textbf{Hubble tuning fork diagram}, which organizes galaxies into distinct families based on their visual properties.

Galaxies are broadly divided into:

\begin{itemize}
    \item \textbf{Elliptical galaxies} (E)
    \item \textbf{Lenticular galaxies} (S0)
    \item \textbf{Spiral galaxies} (S)
    \item \textbf{Irregular galaxies} (Irr)
\end{itemize}

\textbf{Elliptical galaxies} are characterized by smooth light distributions and elliptical shapes. Their morphology is quantified through the eccentricity:
$$
E = 1 - \frac{b}{a}
$$
where $a$ and $b$ are the major and minor axes. The classification ranges from nearly spherical (E0) to highly elongated (E7).

\textbf{Spiral galaxies} exhibit a disk structure with spiral arms and are further subdivided into:
\begin{itemize}
    \item Early spirals: Sa, Sb
    \item Late spirals: Sc, Sd
\end{itemize}

A subclass of spirals, the \textbf{barred spirals} (SB), shows a central bar structure from which the arms originate.

\textbf{Irregular galaxies} lack a well-defined structure and are often classified into Irr I and Irr II depending on their degree of irregularity.

\begin{observationblock}[Projection effects]
The observed morphology of a galaxy depends on its orientation with respect to the line of sight. A disk galaxy can appear circular (face-on) or highly flattened (edge-on), which must be considered when interpreting observations.
\end{observationblock}

\section{Color Classification}

In addition to morphology, galaxies can be classified according to their \textbf{color}, which reflects the properties of their stellar populations.

A commonly used color index is $B - V$, which measures the difference between blue and visible magnitudes. This quantity provides insight into the dominant stellar population:

\begin{itemize}
    \item \textbf{Red galaxies}: dominated by older, cooler stars
    \item \textbf{Blue galaxies}: dominated by young, hot stars
\end{itemize}

Early-type galaxies (ellipticals and lenticulars) typically exhibit red colors, indicating old stellar populations and little ongoing star formation. In contrast, late-type galaxies (spirals and irregulars) are generally bluer due to active star formation in their disks.

\begin{observationblock}[Dust effects]
The presence of dust can significantly affect the observed color of a galaxy. Dust absorbs and scatters blue light more efficiently, making galaxies appear redder than their intrinsic color.
\end{observationblock}

\missing{Lecture 21/10}

\section{Luminosity Function of Galaxies}

The \textbf{luminosity function} $\phi(L)$ describes the number density of galaxies as a function of their luminosity. It is a fundamental statistical tool that connects observations with galaxy formation models.

Formally, $\phi(L)\,dL$ represents the number density of galaxies with luminosities between $L$ and $L + dL$. The total number density is obtained by integrating over all luminosities:
$$
n = \int_0^{\infty} \phi(L)\, dL
$$

Similarly, the total luminosity density is:
$$
\rho_L = \int_0^{\infty} L\, \phi(L)\, dL
$$

A widely used parametrization is the \textbf{Schechter luminosity function}:
$$
\phi(L) = \phi^* \left(\frac{L}{L^*}\right)^{\alpha} \exp\left(-\frac{L}{L^*}\right)
$$

Here:
\begin{itemize}
    \item $\phi^*$ is a normalization constant
    \item $L^*$ is a characteristic luminosity marking the transition between power-law and exponential behavior
    \item $\alpha$ determines the slope at low luminosities
\end{itemize}

\begin{observationblock}[Physical interpretation]
The Schechter function reflects two regimes: a power-law increase at low luminosities, dominated by numerous faint galaxies, and an exponential cutoff at high luminosities, indicating the rarity of very bright systems.
\end{observationblock}

Determining the luminosity function observationally involves several challenges, including distance uncertainties, selection effects, and the need to probe sufficiently large volumes to account for large-scale structure.

\section{Systems of Galaxies}

Galaxies are not randomly distributed in space but instead form \textbf{gravitationally bound systems}.

Two main categories are identified:

\subsection*{Groups of galaxies}

Groups are relatively small systems characterized by:
\begin{itemize}
    \item Number of galaxies: $N \lesssim 50$
    \item Typical masses: $M \gtrsim 10^{13} M_\odot$
    \item Velocity dispersions: $\sigma_v \sim 200\text{--}300\ \text{km s}^{-1}$
    \item Temperatures: $T \sim 1\text{--}3\ \text{keV}$
\end{itemize}

\subsection*{Clusters of galaxies}

Clusters are much larger and more massive systems:
\begin{itemize}
    \item Number of galaxies: $N \gtrsim 50$
    \item Typical masses: $M \gtrsim 10^{14} M_\odot$
    \item Velocity dispersions: $\sigma_v \sim 400\text{--}1000\ \text{km s}^{-1}$
    \item Temperatures: $T \sim 4\text{--}10\ \text{keV}$
\end{itemize}

\begin{observationblock}[Multi-component nature]
Galaxy systems are composed of multiple components: stars, hot intracluster gas, and dark matter. The latter dominates the total mass budget, accounting for the majority of the gravitational potential.
\end{observationblock}

These systems provide key laboratories for studying large-scale structure formation and the interplay between baryonic and dark matter components.

% stop here